\subsection{矩阵的迹}

设 C 为一个交换环；对于每一个方阵 \(X = (x_{ij})\) （C 上的，对应于有限指标集 I），X 的迹是元素

\[
\operatorname{Tr} (X) = \sum_ {i \in I} x _ {i i}.\tag{49}
\]

设 E 是一个具有有限基 \((e_{i})_{i\in I}\) 的 C-模；对 E 的每一个自同态 u，

\[
\operatorname{Tr} (u) = \operatorname{Tr} (M (u))\tag{50}
\]

\(M(u)\) 是 u 关于基 \((e_{i})\) 的矩阵；这由 §4 第 3 号公式 (17) 立即得出，只要把该公式应用于自同态 \(x \mapsto \langle x, e_{i}^{*} \rangle e_{j}\) （其中 \(e_{i}^{*}\) 是 \((e_{i})\) 的对偶基）；由此借助线性性过渡到一般情形。公式 (49) 表明

\[
\operatorname{Tr} (u) = \sum_ {i} \left\langle u (e _ {i}), e _ {i} ^ {*} \right\rangle\tag{51}
\]

对 E 的每一个基 \((e_i)\) 成立（参见 §4，第 3 号，公式 (17)）。

若 X 是 C 上 \((m, n)\) 型的矩阵，且 Y 是 C 上 \((n, m)\) 型的矩阵，则

\[
\operatorname{Tr} (X Y) = \operatorname{Tr} (Y X)\tag{52}
\]

由上述事实及 §4 第 3 号的命题 3 可知；(52) 也可直接得出，因为若 \(X = (x_{ij})\) , \(Y = (y_{ji})\) ( \(1 \leqslant i \leqslant m\) , \(1 \leqslant j \leqslant n\) )，则

\[
\operatorname{Tr} (X Y) = \sum_ {i, j} x _ {i j} y _ {j i}\tag{53}
\]

由 (49) 即得。后一公式还证明了：

命题7. 设 C 是一个交换环，对每个矩阵 \(P \in \mathbf{M}_n(\mathbf{C})\) ，令 \(f_P\) 为线性形式 \(X \mapsto \operatorname{Tr}(PX)\) （在 \(\mathbf{M}_n(\mathbf{C})\) 上）；映射 \(P \mapsto f_P\) 是 \(\mathbf{M}_n(\mathbf{C})\) 到其对偶上的 C-线性双射。

命题8. 若 \(g\) 是 C-模 \(\mathbf{M}_n(\mathbf{C})\) 上的一个线性形式，使得 \(g(XY) = g(YX)\) 对所有矩阵 \(X, Y\) （属于 \(\mathbf{M}_n(\mathbf{C})\) ）成立，则存在且仅存在一个标量 \(c \in \mathbf{C}\) ，使得 \(g(X) = c\) 。Tr(X) 对每个矩阵 \(X \in \mathbf{M}_n(\mathbf{C})\) 成立。

由于该命题在 \(n = 1\) 时是平凡的，故可只限于 \(n \geqslant 2\) 的情形。取 \(X = E_{ij}\) , \(Y = E_{jk}\) （其中 \(i \neq k\) ），得到 \(g(E_{ik}) = 0\) ；再取 \(X = E_{ij}\) , \(Y = E_{ji}\) （其中 \(i \neq j\) ），得到 \(g(E_{ii}) = g(E_{jj})\) ；由于诸 \(E_{ij}\) 构成 \(\mathbf{M}_n(\mathbf{C})\) 的一组基，命题立即得证。

